\scanpage{201}
\mspage{10.}
$C_{S'}$ défini ainsi :
\[
\begin{aligned}
\mathrm{Ob}(C_{S'})&=\{\underline a=(a_i)_{i\in\unsure{I|J}}\mid \forall i,\ a_i\ \text{une équation de }D_i\}\\
\mathrm{Hom}(\underline a,\underline b)&=\{\unsure{\zeta|\xi}=(\zeta_i)_{i\in\unsure{J|I}},\mid\forall i,\ b_i=a_i\zeta_i^{n_i}\}
\end{aligned}
\]

\deleted{Image inverse et} Composition des hom évidente. Images inverses aussi. On trouve comme champ associé une gerbe abélienne de groupe $\mu_{\underline n}$. (L'existence d'une section \uline{dépend} de la nullité d'un élément de $H^2(X,\mu_{\underline n})$, qui n'est autre, comme on devine, que la famille des classes de cohomologie des $D_i$ dans les $H^{\unsure{2|1}}(X,\mu_{n_i})$ !) : \uline{donnée de ramification kummérienne} définie par $(D_i)_{i\in I}=D$ et $\underline n$.

\uline{Ex 2 bis} On donne $D=(D_i)_{i\in I}$, on passe à la limite sur $\underline n$ : donnée de ramification kummérienne définie par $D$.

Ces exemples boitent un tout petit peu, on voudrait en fait que $n_i$ soit premier aux car résiduelles de $D_i$ seulement, pas néc. celles de tt $S$. Dans le cas de l'exemple 1, on prendra les rev. étales de $Z_{\underline n}^{\underline a}$, sur lesquels $\mu_n$ opère de façon compatible avec ses opérations sur $Z_{\underline n}^{\underline a}$. Dans le cas de l'exemple 2
% archive sheet number: 200
